% Exam Template for UMTYMP and Math Department courses
%
% Using Philip Hirschhorn's exam.cls: http://www-math.mit.edu/~psh/#ExamCls
%
% run pdflatex on a finished exam at least three times to do the grading table on front page.
%
%%%%%%%%%%%%%%%%%%%%%%%%%%%%%%%%%%%%%%%%%%%%%%%%%%%%%%%%%%%%%%%%%%%%%%%%%%%%%%%%%%%%%%%%%%%%%%

% These lines can probably stay unchanged, although you can remove the last
% two packages if you're not making pictures with tikz.
\documentclass[11pt]{exam}
\RequirePackage{amssymb, amsfonts, amsmath, latexsym, verbatim, xspace, setspace}
\RequirePackage{tikz, pgflibraryplotmarks}

% By default LaTeX uses large margins.  This doesn't work well on exams; problems
% end up in the "middle" of the page, reducing the amount of space for students
% to work on them.
\usepackage[margin=1in]{geometry}
\usepackage{enumerate}
\usepackage{amsthm}

\theoremstyle{definition}
\newtheorem{soln}{Solution}

% Here's where you edit the Class, Exam, Date, etc.
\newcommand{\class}{Math 407 Section 1}
\newcommand{\term}{Fall 2026}
\newcommand{\examnum}{Exam I}
\newcommand{\examdate}{February 11, 2026}
\newcommand{\timelimit}{1 Hour 50 Minutes}
\newcommand{\ol}[1]{\overline{#1}}

% For an exam, single spacing is most appropriate
\singlespacing
% \onehalfspacing
% \doublespacing

% For an exam, we generally want to turn off paragraph indentation
\parindent 0ex

\begin{document} 

% These commands set up the running header on the top of the exam pages
\pagestyle{head}
\firstpageheader{}{}{}
\runningheader{\class}{\examnum\ - Page \thepage\ of \numpages}{\examdate}
\runningheadrule

\begin{flushright}
\begin{tabular}{p{2.8in} r l}
\textbf{\class} & \textbf{Name (Print):} & \makebox[2in]{\hrulefill}\\
\textbf{\term} &&\\
\textbf{\examnum} & \textbf{Student ID:}&\makebox[2in]{\hrulefill}\\
\textbf{\examdate} &&\\
\textbf{Time Limit: \timelimit} % & Teaching Assistant & \makebox[2in]{\hrulefill}
\end{tabular}\\
\end{flushright}
\rule[1ex]{\textwidth}{.1pt}


This exam contains \numpages\ pages (including this cover page) and
\numquestions\ problems.  Check to see if any pages are missing.  Enter
all requested information on the top of this page, and put your initials
on the top of every page, in case the pages become separated.\\

You may \textit{not} use your books or notes on this exam.\\

You are required to show your work on each problem on this exam.  The following rules apply:\\

\begin{minipage}[t]{3.7in}
\vspace{0pt}
\begin{itemize}

%\item \textbf{If you use a ``fundamental theorem'' you must indicate this} and explain
%why the theorem may be applied.

\item \textbf{Organize your work}, in a reasonably neat and coherent way, in
the space provided. Work scattered all over the page without a clear ordering will 
receive very little credit.  

\item \textbf{Mysterious or unsupported answers will not receive full
credit}.  A correct answer, unsupported by calculations, explanation,
or algebraic work will receive no credit; an incorrect answer supported
by substantially correct calculations and explanations might still receive
partial credit.  This especially applies to limit calculations.

\item If you need more space, use the back of the pages; clearly indicate when you have done this.

\item If the problem asks for a proof, be sure to carefully justify your work, including any theorems from class.

Note: you may NOT use a theorem or result from class to prove something when it makes the problem entirely trivial.  If you are unsure whether a particular theorem or result is allowed, just ask!

\end{itemize}

Do not write in the table to the right.
\end{minipage}
\hfill
\begin{minipage}[t]{2.3in}
\vspace{0pt}
%\cellwidth{3em}
\gradetablestretch{2}
\vqword{Problem}
\addpoints % required here by exam.cls, even though questions haven't started yet.	
\gradetable[v]%[pages]  % Use [pages] to have grading table by page instead of question

\end{minipage}
\newpage % End of cover page

%%%%%%%%%%%%%%%%%%%%%%%%%%%%%%%%%%%%%%%%%%%%%%%%%%%%%%%%%%%%%%%%%%%%%%%%%%%%%%%%%%%%%
%
% See http://www-math.mit.edu/~psh/#ExamCls for full documentation, but the questions
% below give an idea of how to write questions [with parts] and have the points
% tracked automatically on the cover page.
%
%
%%%%%%%%%%%%%%%%%%%%%%%%%%%%%%%%%%%%%%%%%%%%%%%%%%%%%%%%%%%%%%%%%%%%%%%%%%%%%%%%%%%%%

\begin{questions}

\addpoints

\question[10]\mbox{}
\textbf{TRUE or FALSE!}  Write  TRUE if the statement is true.  Otherwise, write FALSE.  Your response should be in ALL CAPS.  No justification is required.
\begin{enumerate}[(a)]
\item  The group $\mathbb Q$ with binary operation $+$ is a cyclic group
\vspace{1.3in}
\item  There are exactly two primitive $4$'th roots of unity
\vspace{1.3in}
\item  A cyclic group must be Abelian
\vspace{1.3in}
\item  $xy+xz+yz$ is an elementary symmetric polynomial in three variables
\vspace{1.3in}
\item  The order of the permutation $(1\ 3\ 4\ 5\ 2)$ is $5$
\vspace{1.3in}
\end{enumerate}

\newpage
\question[10]\mbox{}
\begin{enumerate}[(a)]
\item  Write the definition of a polynomial $p(z_1,\dots, z_n)$ being symmetric
\vspace{1in}
\item  State Vieta's Theorem
\vspace{3in}
\item  Suppose that $z_1,z_2,z_3,z_4$ are the solutions of
$$z^4 + 3z^3 + 2z^2 + 5z + 7 = 0.$$
Determine the value of
$$z_1^2+z_2^2+z_3^2+z_4^2$$
\end{enumerate}

\newpage
\question[10]\mbox{}
Let $G$ be a group and suppose that $\varnothing\neq H\subseteq G$ satisfies $ab^{-1}\in H$ for all $a,b\in H$.  Prove that $H$ is a subgroup of $G$.

\newpage
\question[10]\mbox{}
Consider the set
$$G = (\mathbb{Z}_{7})^\times = \{x\in\mathbb{Z}_{7}: x\neq 0\}.$$
with the binary operation given by multiplication modulo $7$.
\begin{enumerate}[(a)]
\item Show that $G$ has an identity element
\vspace{1.5in}
\item For each element of $G$, determine the inverse.
\vspace{3in}
\item Prove that $G$ is a cyclic group
\vspace{2in}
\item Explain why $G$ is isomorphic to $\mathbb{Z}_6$
\end{enumerate}

\newpage
\question[10]\mbox{}
\begin{enumerate}[(a)]
\item Write down the definition of an $n$'th root of unity and a primitive $n$'th root of unity.
\vspace{2in}
\item Write down all $15$'th roots of unity.  How many are there?
\vspace{3in}
\item Which of the $15$'th roots of unity are primitive?  How many are there?
\end{enumerate}

\newpage
\question[10]\mbox{}
Consider the group $G=S_9$.
\begin{enumerate}[(a)]
\item Rewrite the following permutation in cycle notation
$$\binom{1\ 2\ 3\ 4\ 5\ 6\ 7\ 8\ 9}{8\ 5\ 7\ 6\ 3\ 2\ 1\ 9\ 4}$$
\vspace{2in}
\item Rewrite the following permutation in standard notation
$$(2\ 4\ 6)(1\ 3)(5\ 7\ 9)$$
\vspace{2in}
\item How many elements are in the cyclic subgroup $\langle \sigma\rangle$ generated by the permutation
$$\sigma = (1\ 2\ 3)(4\ 5)$$
\end{enumerate}

\newpage
\question[10]\mbox{}
\begin{enumerate}[(a)]
\item Write down the definition of the sign of a permutation in $S_n$.
\vspace{1in}
\item Prove that the following set is a subgroup of $S_n$
$$A_n = \{\sigma\in S_n:\ \text{$\sigma$ is even}\}$$
\vspace{3in}
\item Prove that for all integers $n > 1$,
$$|A_n| = \frac{n!}{2}$$
\end{enumerate}

\end{questions}
\end{document}
